\documentclass[11pt,letterpaper]{article}
\usepackage[T1]{fontenc}
\usepackage{lmodern}
\usepackage[margin=0.94in,headheight=15pt]{geometry}
\usepackage{microtype}
\usepackage{amsmath,amssymb,amsthm,mathtools}
\usepackage{booktabs,longtable,tabularx,array}
\usepackage{enumitem}
\usepackage[dvipsnames]{xcolor}
\usepackage{listings}
\usepackage{xurl}
\usepackage{hyperref}
\usepackage{fancyhdr}
\hypersetup{colorlinks=true,linkcolor=MidnightBlue,citecolor=MidnightBlue,
 urlcolor=MidnightBlue,pdftitle={Constructive Mordell-Weil descent in iPRA},
 pdfauthor={}}
\pagestyle{fancy}
\fancyhf{}
\fancyhead[L]{\small Constructive Mordell--Weil descent}
\fancyhead[R]{\small iPRA}
\fancyfoot[C]{\thepage}
\setlength{\parskip}{0.35em}
\setlength{\parindent}{1.1em}
\setlist{itemsep=0.2em,topsep=0.3em}
\newcommand{\N}{\mathbb N}
\newcommand{\Z}{\mathbb Z}
\newcommand{\Q}{\mathbb Q}
\newcommand{\iPRA}{\mathsf{iPRA}}
\newcommand{\PR}{\mathsf{PR}}
\newcommand{\dotminus}{\mathbin{\mathord{\dot-}}}
\newcommand{\norm}[1]{\left\lVert#1\right\rVert}
\newcommand{\abs}[1]{\left|#1\right|}
\newcommand{\code}[1]{\texttt{\small #1}}
\DeclareMathOperator{\ev}{eval}
\DeclareMathOperator{\htwo}{ceilLog_2}
\DeclareMathOperator{\bl}{bitlen}
\DeclareMathOperator{\Cover}{Cover}
\DeclareMathOperator{\Good}{Good}
\DeclareMathOperator{\Halves}{Halves}
\DeclareMathOperator{\Sha}{Sha}
\DeclareMathOperator{\adj}{adj}
\newtheorem{theorem}{Theorem}[section]
\newtheorem{lemma}[theorem]{Lemma}
\newtheorem{corollary}[theorem]{Corollary}
\newtheorem{proposition}[theorem]{Proposition}
\theoremstyle{definition}
\newtheorem{definition}[theorem]{Definition}
\newtheorem{remark}[theorem]{Remark}
\lstset{basicstyle=\ttfamily\footnotesize,breaklines=true,
 columns=fullflexible,keepspaces=true,frame=single,rulecolor=\color{gray!40},
 showstringspaces=false}

\title{\textbf{Constructive Mordell--Weil descent\\
 in intuitionistic primitive recursive arithmetic}\\[0.6em]
 \large Explicit generation from quotient representatives,\\
 and an unconditional family of rational Mordell--Weil groups}
\author{}
\date{November 15, 2026}

\begin{document}
\maketitle
\begin{abstract}
We give a precise, witness-producing version of the height-descent part of
Mordell--Weil, with an arithmetic proof using primitive recursion and
quantifier-free induction. For every nonsingular integral short Weierstrass
equation, we construct explicit integer height constants, a complete finite
algorithm for rational division by two, and an algorithm which turns a
supplied covering list modulo doubling into generators and coordinates of
every rational point. The universal covering assertion remains an explicit
hypothesis. In particular, no selector or height estimate is hidden in that
hypothesis.

Independently, we prove in full that the rational group of
$y^2=x^3-x$ is $(\Z/2\Z)^2$, by a bounded version of Fermat's square-area
descent. This gives unconditional uniform classifications for rationally
scaled models and their finite products, in every finite dimension.
We give precise coding conventions, induction invariants, executable exact
arithmetic, a compiled Lean descent kernel, and a module-by-module plan for
the deeply embedded iPRA formalization in ProveIt.

This is not an unconditional constructive proof of finite generation for
every abelian variety over every number field. Producing the quotient
representatives uniformly is not solved here. The distinction is necessary:
a general algorithm producing Mordell--Weil generators is still an open
effectivity problem already for elliptic curves over $\Q$.
\end{abstract}

\begin{center}
\fbox{\begin{minipage}{0.91\textwidth}
\textbf{Verification boundary.}
The arithmetic arguments in this paper are human-readable proofs with
explicit iPRA induction invariants. The accompanying Lean files check
selected native Lean theorems. The complete arithmetic results have not
yet been assembled as values of ProveIt's embedded \code{Prf PRAAx}.
No independence or impossibility result about iPRA is asserted.
\end{minipage}}
\end{center}

\clearpage
\begingroup
\small
\setlength{\parskip}{0pt}
\tableofcontents
\endgroup
\newpage
\section{The theorem being proved and its logical scope}
\label{sec:scope}

The classical Mordell--Weil theorem says that $A(K)$ is a finitely generated
abelian group whenever $A$ is an abelian variety over a number field $K$.
The requested Wikipedia page supplies this formulation
\cite{wikipedia}; Milne gives a detailed treatment of the elliptic case
\cite{milne}. An arithmetic constructive interpretation should say what
finite data are produced, how an arbitrary rational point is expressed in
those data, and why every search terminates within a primitive recursive
bound.

For an explicitly coded group, a strong form is a finite list
$G=(G_0,\ldots,G_{s-1})$ and a total primitive recursive function
$\operatorname{coord}$ satisfying
\[
 P=\sum_{j<s}\operatorname{coord}(P)_jG_j
 \qquad\text{for every valid point code }P.
 \tag{1.1}
\]
The coefficients may be signed integers. Our descent actually produces
nonnegative coefficients; negative points are allowed among the generators.
An equivalent form useful for formalization returns a finite expression
formed by addition and doubling, whose leaves belong to $G$.

\subsection{Why the unrestricted constructive theorem is not claimed}

The ordinary proof has two major ingredients: finiteness of
$A(K)/mA(K)$, and height descent. Passing from an upper bound on the number
of cosets to a list of \emph{actual rational representatives} is a real
constructive issue. Sutherland explicitly identifies the absence of an
effective representative-height bound in the usual argument
\cite{sutherland}. The current official Magma handbook describes the
general algorithmic Mordell--Weil problem as open \cite{magma}.

These observations justify a boundary, not an unprovability claim.
In particular, we do not infer that the classical statement is independent
of iPRA, that every possible constructive reformulation is impossible, or
that an arbitrary existential theorem of intuitionistic arithmetic has an
immediate primitive recursive witness-extraction theorem. Our claim is
more direct: the proposed unrestricted \emph{algorithmic} conclusion would
itself supply a currently unavailable uniform generator algorithm.

The elementary logical distinction can be seen without elliptic curves.
Let $f_e(t)$ be $1$ if a fixed machine with code $e$ halts within $t$
steps, and $0$ otherwise. This is a primitive recursive function of $e,t$,
and its image has at most two elements. A uniform procedure returning
finitely many times whose values cover the whole image would decide
halting: simulate to the greatest returned time. Thus an explicitly finite
ambient set does not by itself give a terminating procedure for listing
which of its labels are attained.

In descent, an explicitly computed finite Selmer group is an ambient
container for rational-point classes; it need not consist entirely of
classes with a rational point. The familiar exact sequence
\[
0\longrightarrow E(\Q)/2E(\Q)
 \longrightarrow \operatorname{Sel}^{(2)}(E/\Q)
 \longrightarrow \Sha(E/\Q)[2]\longrightarrow0
\]
expresses that distinction in the classical theory \cite{milne}.
We use this only to explain the missing input. No cohomology theorem is
used as an unproved arithmetic premise below.

\subsection{The three completed mathematical results}

\begin{enumerate}
\item \textbf{A general arithmetic descent schema.}
For any fixed primitive recursively presented abelian group, a finite
bounded-height enumeration and a verified contracting decomposition give
generators, coefficients, and explicit descent fuel
(Theorem~\ref{thm:abstract}).

\item \textbf{A supplied-cover theorem for every integral elliptic model.}
For $E:y^2=x^3+Ax+B$, $4A^3+27B^2\ne0$, the only nonautomatic arithmetic
input is a finite list $R$ of rational points satisfying
\[
 \Cover_E(R):\quad
 \forall P\in E(\Q)\ \exists i<|R|\ \exists Q\in E(\Q)\
       P=2Q+R_i.
 \tag{1.2}\label{eq:cover}
\]
Explicit polynomial formulas compute the needed height constants and
a primitive recursive selector. We then compute a finite generating list
and coordinates (Theorem~\ref{thm:elliptic-cover}). The condition
\eqref{eq:cover} is an explicit universal hypothesis, not a claim that a
finite test decides it.

\item \textbf{Unconditional arithmetic Mordell--Weil instances.}
We prove
\[
 E_1(\Q)=\{O,(0,0),(1,0),(-1,0)\},\qquad E_1:y^2=x^3-x,
\]
and identify the group as $(\Z/2\Z)^2$. For any finite list
$u_1,\ldots,u_r\in\Q^\times$, we explicitly classify the rational group
of $\prod_{j=1}^r E_{u_j}$, where
$E_u:y^2=x^3-u^4x$, as $(\Z/2\Z)^{2r}$
(Theorem~\ref{thm:products}).
\end{enumerate}

All arithmetic dependencies needed for these statements are explained
below. General number fields and arbitrary abelian varieties fit the
abstract interface only after their corresponding arithmetic data and
proofs have been supplied. We do not claim to construct those data
uniformly.

\section{The arithmetic setting and coding conventions}
\label{sec:foundations}

\subsection{Logic and nonlogical axioms}

Our object theory is the first-order intuitionistic primitive recursive
arithmetic implemented in ProveIt's \code{ConciseIPRA/Lean/ConciseIPRA/PR.lean}
\cite{proveit}. Its language has equality and a function symbol for each
primitive recursive definition. Symbols are generated from zero,
successor and projections by composition and primitive recursion:
\[
\begin{split}
 \operatorname{Rec}(f,g)(0,\vec x)&=f(\vec x),\\
 \operatorname{Rec}(f,g)(n+1,\vec x)
   &=g(n,\operatorname{Rec}(f,g)(n,\vec x),\vec x).
\end{split}
\]
The nonlogical axioms are these defining equations, successor injectivity,
$Sx\ne0$, and induction for quantifier-free formulas, with arbitrary
numerical parameters:
\[
 \varphi(0)\ \to\
 \bigl(\forall n\,(\varphi(n)\to\varphi(n+1))\bigr)
 \ \to\ \forall n\,\varphi(n).
 \tag{2.1}\label{eq:qfind}
\]
The logic is intuitionistic first-order natural deduction with equality,
$\bot,\wedge,\vee,\to,\forall,\exists$, the usual introduction and
elimination rules, and falsity elimination. No excluded-middle schema,
choice schema, Markov principle, or double-negation elimination is added.

Numerical comparisons and finite tests are decidable. Their characteristic
functions are primitive recursive; bounded conjunctions and disjunctions
are finite folds. When a displayed induction invariant uses a bounded
quantifier, it is shorthand for the corresponding Boolean test equalling
$1$. This conversion is part of the proof plan, not permission to use
unrestricted induction.

\subsection{Exact codes and total operations}

We fix the following concrete conventions; alternative primitive recursive
encodings would give the same results.

\begin{itemize}
\item Encode an integer $z$ by $2z$ for $z\ge0$, and $2|z|-1$ for $z<0$.
Sign, magnitude, addition and multiplication are primitive recursive.
\item Encode a pair by
$\langle a,b\rangle=2^a(2b+1)-1$. Recover $a$ by bounded division of the
successive powers of $2$ into the positive number $n+1$, then recover $b$.
The inverse searches can use $n+1$ as fuel.
\item The empty list has code $0$ and
$\operatorname{cons}(a,\ell)=1+\langle a,\ell\rangle$.
The tail of a nonempty list has smaller code. Iterating at most the input
code therefore suffices for length, access, folds and validity checks.
\item A rational is a reduced pair $(a,b)$ with $a\in\Z$, $b>0$ and
$\gcd(|a|,b)=1$. Normalize by a positive gcd. Equality is equality of
canonical pairs; addition and multiplication use cross products followed
by normalization.
\item An elliptic point is the tag $O$, or an affine pair of canonical
rationals satisfying the curve equation. We may use code $0$ for $O$
and $1+\langle x,y\rangle$ for an affine pair. Invalid codes carry a
failure flag and return specified default outputs.
\end{itemize}

Each program is a fixed primitive recursive definition. Statements about
an arbitrary primitive recursively presented group are schemas, one for
each fixed collection of programs and proved laws. We do not postulate a
single total primitive recursive evaluator for all primitive recursive
program codes. For the elliptic theorem the programs are fixed, with
numeric $A,B,R,P$ as ordinary inputs.

Finite list construction, an exponential-size search, or a very large
integer is compatible with primitive recursion. No claim of practical
running time or polynomial complexity is intended.

\subsection{Bounded arithmetic lemmas}
\label{subsec:gcd}

Here are the elementary ingredients used repeatedly later. Giving their
witnesses is useful both mathematically and for the object formalization.

\begin{lemma}[Gcd with Bezout witnesses]\label{lem:gcd}
For nonnegative $a,b$, not both zero, one computes $g>0$ and signed
$u,v$ with
\[
g\mid a,\quad g\mid b,\quad ua+vb=g,
\]
and every common divisor divides $g$. The computation and proof are
primitive recursive and use quantifier-free induction.
\end{lemma}
\begin{proof}
Handle a zero argument directly. Otherwise order the entries as
$(\max(a,b),\min(a,b))$ and run the Euclidean step
$(a_i,b_i)\mapsto(b_i,a_i\bmod b_i)$, freezing at $b_i=0$.
When active, the positive sum of the entries strictly decreases, so
$a+b+1$ stages suffice. Track both entries as integer linear combinations
of the original arguments. If $a_i=q_ib_i+r_i$, replace the second pair
of coefficients by the old first pair minus $q_i$ times the old second
pair. The linear-combination equations are preserved.

For a fixed numerical parameter $d$, the divisibility of both entries by
$d$ is preserved in both directions. At the terminal pair $(g,0)$ this
proves the gcd property. The invariant
$b_i=0\ \vee\ a_i+b_i+i\le a+b$, together with the two tracked equations
and Boolean divisibility predicates, is quantifier-free. Division itself
is a scan for a quotient at most the dividend.
\end{proof}

\begin{lemma}[Coprime cancellation and square splitting]\label{lem:squares}
The following operations have primitive recursive witnesses:
\begin{enumerate}
\item If $\gcd(a,b)=1$ and $a\mid bc$, then $a\mid c$.
\item If $d>0$ and $d^2\mid c^2$, then $d\mid c$.
\item If $a,b,c>0$, $\gcd(a,b)=1$ and $ab=c^2$, there are positive
$r,s$ with $a=r^2$, $b=s^2$ and $c=rs$.
\end{enumerate}
\end{lemma}
\begin{proof}
For (1), multiply a Bezout equation for $a,b$ by $c$.
Squaring a Bezout equation also shows that coprimality with $b$ implies
coprimality with $b^2$; multiplying two Bezout equations shows that
coprimality with each of two factors implies coprimality with their product.

For (2), the case $c=0$ is immediate. Otherwise write
$c=hr,d=hs$, where $h=\gcd(c,d)$ and $\gcd(r,s)=1$.
The hypothesis gives $s^2\mid r^2$, so $s\mid r^2$.
But $s$ is coprime to $r^2$; Bezout then gives $s\mid1$, hence $s=1$.

For (3), put $h=\gcd(a,c)$, $a=hu,c=hv$, with $\gcd(u,v)=1$.
The equation gives $ub=hv^2$, so (1) implies $u\mid h$.
Write $h=ue$; cancellation gives
$a=eu^2,b=ev^2,c=euv$. The positive integer $e$ divides both $a,b$,
so $e=1$. Thus $r=u,s=v$ work. All quotients are exact bounded divisions.
These are equation manipulations and applications of Lemma~\ref{lem:gcd};
no factorization theorem or least-counterexample principle is needed.
\end{proof}

The last conclusion applies successively to any fixed number of pairwise
coprime positive factors whose product is a square. An exact integer
square root is also primitive recursive: scan $0,\ldots,N$ for $r^2=N$,
with a separate success bit. A rational square root is found by applying
this to the nonnegative numerator and positive denominator of a reduced
fraction. Lemma~\ref{lem:squares}(2) proves that a rational with integral
square is an integer. This proves completeness of the rational-square
test, not just soundness.

\begin{lemma}[Finite scans]\label{lem:scan}
For a fixed primitive recursive Boolean predicate $T(j,\vec a)$ and an
input bound $B$, scanning $0,\ldots,B$ returns the first successful index
and a success flag. If a successful index exists within that interval,
the flag is $1$ and the returned index satisfies $T$.
\end{lemma}
\begin{proof}
Store a success bit and index, freezing after the first success. The
Boolean invariant states that success is the disjunction of the tests
already inspected, and that a successful stored index is within that
prefix and passes the test. Finite disjunction is itself a primitive
recursive fold; induction on the scan index is quantifier-free.
Existential elimination may then be used outside that induction to show
that a supplied bounded witness forces success.
\end{proof}

\section{A primitive recursive Mordell--Weil descent schema}
\label{sec:descent}

Let $\Gamma$ be an abelian group represented by valid natural-number codes,
with fixed primitive recursive equality, addition, negation and scalar
multiplication. Suppose its laws are proved on valid codes in iPRA.
The group is not required to be finite or to have a decidable rank.
Let $h:\Gamma\to\N$ be a fixed primitive recursive height.

An enumeration $L(B)$ means a primitive recursive finite list of valid
points, with the proved property
\[
 h(P)\le B\quad\Longrightarrow\quad P\text{ occurs in }L(B).
 \tag{3.1}\label{eq:enum}
\]
Duplicates are harmless. Including the identity in every list avoids
irrelevant empty-list branches.

\begin{theorem}[Explicit contracting descent]\label{thm:abstract}
Suppose $R=(R_0,\ldots,R_{r-1})$, $r>0$, is a finite list, $C\in\N$,
and fixed primitive recursive functions $i(P),D(P)$ satisfy, for all
valid $P$ with $h(P)>C$,
\[
 i(P)<r,\qquad D(P)\in\Gamma,\qquad
 P=2D(P)+R_{i(P)},\qquad
 2h(D(P))\le h(P)+C.
 \tag{3.2}\label{eq:contract}
\]
Assume \eqref{eq:enum} and these statements have iPRA proofs.
Then $G=R\mathbin{+\!\!+}L(C)$ generates $\Gamma$. A primitive recursive
function produces the coefficients in \eqref{eq:1-coord} below, with an
iPRA proof of correctness.

More explicitly, write $e=(h(P)\dotminus C)$ and let
$t=\bl(e)$, where $\bl(0)=0$ and, for $e>0$, $\bl(e)$ is the unique
integer with $2^{t-1}\le e<2^t$. The algorithm takes at most $t$ active
descent steps. Every resulting coefficient is at most $2^t$; if $e>0$,
$2^t\le2e$, and if $e=0$ the largest coefficient is $1$.
\end{theorem}

\begin{proof}
Define the excess $e(Q)=h(Q)\dotminus C$. An active step has
$h(Q)=C+e(Q)$, so \eqref{eq:contract} gives
$2h(D(Q))\le2C+e(Q)$. Therefore
\[
 e(D(Q))\le \left\lfloor e(Q)/2\right\rfloor.
 \tag{3.3}\label{eq:halve-excess}
\]
If the child has height below $C$, the left side is zero; otherwise this
is direct subtraction. In particular every active step strictly decreases
height. Freeze the whole state as soon as its height is at most $C$.

The state consists of a current point $Q$, an active-step count $k$,
and an $r$-entry coefficient vector $a$. Start at
$Q=P,k=0,a=0$. An active step, using the old point, performs
\[
 j=i(Q),\quad Q'=D(Q),\quad
 a'_j=a_j+2^k,\quad a'_l=a_l\ (l\ne j),\quad k'=k+1.
 \tag{3.4}\label{eq:3-update}
\]
Run this stopped transition for exactly $t$ stages. The natural number $t$
is computed by a bounded scan: search $0,\ldots,e+1$ for the least $j$
with $e<2^j$. Induction proves $2^{e+1}>e$, so the scan succeeds.

After $j$ stages, \eqref{eq:halve-excess} gives
$e(Q)\le\lfloor e/2^j\rfloor$. A convenient entirely integral invariant
is $2^j e(Q)\le e$, maintained also through frozen stages because then
$e(Q)=0$. At $j=t$, $e<2^t$, so $e(Q)=0$ and $Q$ belongs to $L(C)$.
If one prefers a simpler bound, $h(P)$ stages suffice by strict decrease.
The compiled Lean kernel uses this coarser bound.

The simultaneous group invariant is
\[
 P=2^kQ+\sum_{l<r}a_lR_l.
 \tag{3.5}\label{eq:group-invariant}
\]
It holds initially. Substituting $Q=2Q'+R_j$ proves it after
\eqref{eq:3-update}; a frozen step leaves it unchanged.
At the end, scan $L(C)$ for $Q$ using decidable equality and
Lemma~\ref{lem:scan}. If $Q=L(C)_b$, then
\[
 P=\sum_{l<r}a_lR_l+2^kL(C)_b.
 \tag{3.6}\label{eq:1-coord}
\]
The digit weights added in active steps are $1,2,\ldots,2^{k-1}$.
Consequently $\sum_l a_l=2^k-1$, and $k\le t$ proves the coefficient
bound. The elementary inequality $2^t\le2e$ for $e>0$ follows from the
minimality of $t$.

For the iPRA proof, encode the vector and state. Validity, the excess
inequality, the coefficient-sum identity and
\eqref{eq:group-invariant} are quantifier-free formulas about the coded
functions. The finite sum is a primitive recursive fold. Correctness of
each update supplies the induction step of \eqref{eq:qfind}. The possibly
quantified global premises remain hypotheses of the derivation; they
are not placed inside its induction formula.
\end{proof}

\subsection{A formulation requiring no group identities}

The core also works for a set with a binary operation $\oplus$, without
associativity or commutativity. Assume a child satisfies
$P=(D(P)\oplus D(P))\oplus R_{i(P)}$.
Construct a word by returning the corresponding bounded-height atom when
small, and otherwise returning
\[
 \operatorname{add}\bigl(\operatorname{double}(w(D(P))),
                         \operatorname{rep}(i(P))\bigr).
 \tag{3.7}
\]
With $h(P)$ fuel this gives an expression evaluating exactly to $P$.
Each active stage adds three constructors, so the compact expression has
at most $3h(P)+1$ constructors. A dedicated \emph{double} constructor
avoids physically copying the child expression.

This is the theorem proved in \code{lean/DescentCore.lean}. In an abelian
group, folding the word into coefficients gives (3.6). For elliptic curves,
the word version isolates the arithmetic chord-and-tangent construction
from the independent proof of all abstract group identities.

\subsection{What counts as a certificate}

A finite list, a constant and program definitions do not make their
universal correctness automatically decidable. The assumptions in
Theorem~\ref{thm:abstract} are mathematical assertions with proofs.
There are two honest uses of the theorem: prove these assertions for a
particular input, or state a conditional theorem with those assertions as
premises.

One may also supply finite formal derivations of the premises. Checking
such a derivation is a syntactic operation, and combining it with a fixed
proof template can produce a new derivation. This does not justify an
internal principle saying that every accepted proof has a true conclusion.
Such a principle would be a reflection assertion and is not used here.

\subsection{A simple infinite test group}

For $\Gamma=\Z^r$, take $h(z)=\sum_j z_j^2$, digits
$(z_j\bmod2)_j\in\{0,1\}^r$, and child
$q_j=\lfloor z_j/2\rfloor$. For an even coordinate,
$2q_j^2\le z_j^2+1$; for $z_j=2q_j+1$ the same inequality follows from
$(2q_j+1)^2+1-2q_j^2=2(q_j+1)^2\ge0$. Summing gives the contract
with $C=r$. The finite cube $[-r,r]^r$, filtered by $\sum_jz_j^2\le r$,
enumerates the base region. This illustrates an infinite group where the
algorithm's hypotheses are elementary and the output is not vacuous.
It is an algorithmic test case, not an elliptic-curve rank computation.

\section{Explicit arithmetic on an arbitrary integral elliptic model}
\label{sec:elliptic}

Throughout this section fix integers $A,B$ with
$D_0=4A^3+27B^2\ne0$, and put
\[
 E:\quad y^2=x^3+Ax+B,\qquad S=1+|A|+|B|.
 \tag{4.1}\label{eq:4-model}
\]
Its projective equation is $Y^2Z=X^3+AXZ^2+BZ^3$.
The only point with $Z=0$ is $O=[0:1:0]$. Every other point has unique
affine coordinates, so the tagged rational codes of
Section~\ref{sec:foundations} cover all rational projective points.

\subsection{The operations and the inverse-translation identity}

Define $-O=O$ and $-(x,y)=(x,-y)$; make $O$ the identity.
For affine $P=(x,y),R=(u,v)$, if $x=u,y=-v$, put $P+R=O$.
Otherwise use
\[
 \lambda=
 \begin{cases}
  (v-y)/(u-x),&x\ne u,\\
  (3x^2+A)/(2y),&P=R,\ y\ne0,
 \end{cases}\qquad
 x'=\lambda^2-x-u,\quad y'=\lambda(x-x')-y.
 \tag{4.2}\label{eq:add}
\]
Set $P+R=(x',y')$. If $x=u$, the two curve equations imply
$(y-v)(y+v)=0$; decidable rational equality therefore leaves exactly the
listed branches. Every denominator is nonzero in its own branch.

Substitute the line $Y=\lambda(X-x)+y$ into the cubic. Its polynomial
has $x,u$ as roots, with multiplicity two in the tangent case, and the
sum of its three roots is $\lambda^2$. Polynomial division therefore
gives the third root $\lambda^2-x-u$. Reflecting its ordinate gives
\eqref{eq:add}, which proves closure by direct substitution. This proof
uses only rational arithmetic and monic division by linear factors.

We will need
\[
 (U+R)-R=U.
 \tag{4.3}\label{eq:inverse-translation}
\]
This identity has a short verification that does not require the full
associativity theorem. If the line through $U,R$ meets the cubic a third
time at $-(U+R)$, its reflection meets it at $-U,-R,U+R$.
Applying the same third-intersection rule to $U+R,-R$ therefore yields
$U$. When two of these intersections coincide, the multiplicity gives
the tangent branch. A vertical line gives the inverse branch; an input
$O$ gives the identity branch. Equivalently these are the finite case
splits of \eqref{eq:add}, verified by the same cubic factorization.
Replacing $R$ by $-R$ also gives $(U-R)+R=U$.
Thus the word-producing theorem below has an elementary arithmetic
proof independent of a separate formalization of associativity.

For the usual abelian-group formulation, \eqref{eq:add} is the
elliptic group operation. The general group laws may be reused in a
standard mathematical development; they are explicitly separate
object-theory obligations if one wants the coefficient-vector version
inside iPRA. The compact-word conclusion itself does not need them.
For the unconditional family in Section~\ref{sec:unconditional}, the
entire group law will instead be a verified four-element table.

\subsection{Square and cube denominators}

\begin{lemma}[Canonical affine form]\label{lem:denominators}
Every rational affine point of \eqref{eq:4-model} has a unique expression
\[
 x=a/d^2,\qquad y=b/d^3,\qquad
 a,b\in\Z,\ d>0,\quad \gcd(|a|,d)=1,
 \tag{4.4}\label{eq:denomform}
\]
where
\[
 b^2=a^3+Aa d^4+Bd^6.
 \tag{4.5}\label{eq:clearedcurve}
\]
The conversion is primitive recursive.
\end{lemma}
\begin{proof}
Write $x=a/c$ in reduced form, $c>0$. Then
\[
 x^3+Ax+B=(a^3+Aac^2+Bc^3)/c^3.
\]
The numerator is coprime to $c$, because it is congruent to $a^3$
modulo $c$ and $\gcd(a,c)=1$. It is consequently coprime to $c^3$.
Write $y=b/s$ in reduced form. Both $b^2/s^2$ and the preceding
fraction are reduced. Uniqueness of reduced fractions gives $s^2=c^3$
and $b^2=a^3+Aac^2+Bc^3$.

Since $c^2\mid s^2$, Lemma~\ref{lem:squares} gives $c\mid s$.
Put $d=s/c$. Then $d^2=c$, so $s=d^3$. This proves the equations.
The zero-ordinate branch causes no exception: reduced zero has
denominator $1$, and the argument gives $c=s=d=1$.
Uniqueness follows from uniqueness of the reduced $x$ denominator,
the positive square root, and then the ordinate.
\end{proof}

The uniqueness of reduced fractions used here is elementary: if
$a/c=b/d$ with both denominators positive and both pairs coprime,
$ad=bc$ and coprime cancellation imply $c\mid d$ and $d\mid c$.
Thus $c=d,a=b$. The denominator proof does not invoke valuations,
prime factorization, or algebraic number theory.

\subsection{Integer height and complete bounded-height enumeration}

For a reduced projective rational pair $[a:c]$, set
$H([a:c])=\max(|a|,|c|)$. For an affine elliptic point use its
$x$-coordinate; put $x(O)=[1:0]$. Thus
\[
 H(O)=1,\qquad H(P)=\max(|a|,d^2)
 \quad\hbox{in \eqref{eq:denomform}}.
 \tag{4.6}
\]
All heights are positive integers. Define
\[
 h(P)=\htwo H(P)=\min\{n\le H(P):H(P)\le2^n\}.
 \tag{4.7}\label{eq:binaryheight}
\]
The minimum is a bounded scan. The elementary inequality $N\le2^N$
for $N\ge1$ proves its success. In particular $h(O)=0$, and
$h(P)=0$ whenever $H(P)=1$. No real logarithm occurs in the definition.

\begin{lemma}[An explicit finite point box]\label{lem:ell-enum}
For an integer $M\ge1$, every affine point with $H(P)\le M$ occurs among
the triples
\[
 |a|\le M,\quad 1\le d\le M,\quad d^2\le M,\quad
 |b|\le SM^2,
 \tag{4.8}\label{eq:box}
\]
filtered by $\gcd(|a|,d)=1$ and \eqref{eq:clearedcurve}.
Adding $O$ gives a complete primitive recursive enumeration
$\mathcal L_E(M)$. Its length is at most
\[
 1+(2M+1)M(2SM^2+1).
 \tag{4.9}
\]
\end{lemma}
\begin{proof}
Lemma~\ref{lem:denominators} gives the form and the first three bounds.
The equation implies
$b^2\le(1+|A|+|B|)M^3=SM^3\le S^2M^4$.
Since all compared magnitudes are nonnegative, $|b|\le SM^2$.
The displayed box has the stated number of candidates before filtering.
Conversely every retained triple gives a rational point of the required
height. Each loop has its displayed numerical bound, and all tests are
decidable. Completeness is the preceding inequality, not a conclusion
drawn from an empirical search.
\end{proof}

For Theorem~\ref{thm:abstract} take $L_E(C)=\mathcal L_E(2^C)$.
One may implement the enumeration more efficiently by testing square
roots of the right side of \eqref{eq:clearedcurve}, rather than scanning
every $b$; the same bound proves completeness.

\subsection{A finite polynomial certificate for duplication}
\label{subsec:dup}

The tangent formula in \eqref{eq:add} gives the homogeneous map
\[
 x(2P)=[F(X,Z):G(X,Z)]\quad\text{when }x(P)=[X:Z],
 \tag{4.10}\label{eq:dupmap}
\]
where
\[
\begin{split}
 F(X,Z)&=X^4-2AX^2Z^2-8BXZ^3+A^2Z^4,\\
 G(X,Z)&=4X^3Z+4AXZ^3+4BZ^4.
\end{split}
\tag{4.11}\label{eq:FG}
\]
Indeed $x(2P)=(3x^2+A)^2/(4y^2)-2x$; substitute the curve equation
and homogenize. At $O$ the pair is $[1:0]$.
If $y=0$, the denominator vanishes and the numerator is
$(3x^2+A)^2$, which is nonzero by nonsingularity. Thus the same
projective formula handles every branch.

Construct the integer $8$ by $8$ matrix $M_{A,B}$ of the linear map
\[
 (U,V)\longmapsto UF+VG,
 \tag{4.12}
\]
where $U,V$ are binary cubics, and the output is a binary form of degree
$7$. Use the input order
$X^3,X^2Z,XZ^2,Z^3$ first for $U$, then for $V$, and output order
$X^7,X^6Z,\ldots,Z^7$. The columns are simply these monomials times
$F$ or $G$. Determinant expansion gives
\[
 R_0=\det M_{A,B}=256(4A^3+27B^2)^2\ne0.
 \tag{4.13}\label{eq:resultant}
\]
This is a fixed-size polynomial computation, not an appeal to a
nonconstructive elimination theorem.

For an alternative check of nonvanishing, put $f=X^3+AX+B$ and
$f'=3X^2+A$. Direct expansion proves
\[
 (27B-18AX)f+(6AX^2-9BX+4A^2)f'=4A^3+27B^2.
 \tag{4.14}
\]
Moreover $F(X,1)=f'(X)^2-8Xf(X)$ and $G(X,1)=4f(X)$.
Thus the dehomogenized polynomials are coprime over $\Q$; their
homogeneous versions have no common factor at infinity since
$F(X,0)=X^4$. This also makes the Sylvester matrix invertible by
finite polynomial division and rational linear algebra.

The first and last columns of $\adj M_{A,B}$ supply integral binary
cubics $U_X,V_X,U_Z,V_Z$ such that
\[
 U_XF+V_XG=R_0X^7,\qquad
 U_ZF+V_ZG=R_0Z^7.
 \tag{4.15}\label{eq:bezout7}
\]
All coefficients are computed by fixed-size determinants. If a positive
integer divides $R_0$ and all the coefficients of the four cubics, divide
both identities by it. Denote the resulting nonzero integer by $R_*$.
No sharp normalization is required; a common gcd suffices.
Let
\[
 L=\max\bigl(1,\norm{U_X}_1+\norm{V_X}_1,
                   \norm{U_Z}_1+\norm{V_Z}_1\bigr).
 \tag{4.16}\label{eq:L}
\]
Here $\norm{\cdot}_1$ is the sum of the absolute coefficients, after any
normalization. This is an explicit positive integer depending only on
$A,B$. With
$M=\max(1,2|A|,8|B|,A^2,4,4|A|,4|B|)$, a coarser completely
displayed upper bound is $L\le8\cdot7!\,M^7$.

\begin{lemma}[Duplication lower bound]\label{lem:dup-height}
For every rational point,
\[
                  H(P)^4\le L\,H(2P).
 \tag{4.17}\label{eq:dupheight}
\]
\end{lemma}
\begin{proof}
Write $x(P)=[a:c]$ with $\gcd(|a|,|c|)=1$. Set
\[
 H=\max(|a|,|c|),\qquad T=\max(|F(a,c)|,|G(a,c)|).
\]
The two evaluations cannot both vanish, because
\eqref{eq:bezout7}, with $R_*$ in place of $R_0$, would force
$a=c=0$. Put $g=\gcd(|F(a,c)|,|G(a,c)|)>0$.
The identities show that $g$ divides both $R_*a^7$ and $R_*c^7$.
Since $a^7,c^7$ are coprime, Bezout gives $g\mid R_*$.
Consequently $g\le|R_*|$.

Choose a coordinate of magnitude $H$ and use its corresponding identity.
Every cubic multiplier is at most its coefficient norm times $H^3$,
so
\[
 |R_*|H^7\le L H^3T,\qquad |R_*|H^4\le LT.
 \tag{4.18}
\]
After reduction, $H(2P)=T/g$. Dividing the second inequality by
$g$, and using $|R_*|/g\ge1$, gives \eqref{eq:dupheight}.
All divisions used in the last statement are exact integer divisions;
one can instead multiply through by $g$ throughout.

The gcd argument controls the cancellation in projective coordinates.
Omitting it would not prove a lower height bound.
\end{proof}

\subsection{A translation upper bound with every constant specified}
\label{subsec:translation}

For a nonidentity rational point $R$, define
\[
 T_R=\max\{4S^2H(R)^2,\ H(2R),\ 1\};
 \qquad T_O=1.
 \tag{4.19}\label{eq:TR}
\]
All these quantities are directly computable.

\begin{lemma}[Translation upper bound]\label{lem:translation}
For every pair of rational points $P,R$,
\[
                  H(P-R)\le T_R H(P)^2.
 \tag{4.20}\label{eq:translation}
\]
\end{lemma}
\begin{proof}
For nonidentity $P,R$, use their canonical forms
$P=(a/d^2,b/d^3)$ and $R=(c/e^2,w/e^3)$.
If their abscissae differ, the secant formula simplifies to
\[
 x(P-R)=
 \frac{x(P)^2x(R)+x(P)x(R)^2+
       A(x(P)+x(R))+2B+2y(P)y(R)}
      {(x(P)-x(R))^2}.
 \tag{4.21}
\]
Multiplying numerator and denominator by $d^4e^4$ gives the integer pair
\[
\begin{split}
 N={}&a^2ce^2+ac^2d^2
       +A(ad^2e^4+cd^4e^2)
       +2Bd^4e^4+2bwde,\\
 D={}&(ae^2-cd^2)^2.
\end{split}
\tag{4.22}\label{eq:ND}
\]
Put $H=H(P)$ and $J=H(R)$. The curve equation gives
$b^2\le SH^3$ and $d^2\le H$, whence
$(|b|d)^2\le SH^4\le S^2H^4$, and therefore $|b|d\le SH^2$.
Likewise $|w|e\le SJ^2$.
The first two terms of $N$ each have magnitude at most $H^2J^2$;
the $A$ terms contribute at most $2|A|H^2J^2$;
the $B$ term at most $2|B|H^2J^2$;
and the last term at most $2S^2H^2J^2$.
Thus
\[
 |N|\le2(S+S^2)H^2J^2\le4S^2H^2J^2,\qquad
 |D|\le4H^2J^2.
 \tag{4.23}
\]
Reduction of the pair can only decrease its height. This proves the bound
for distinct abscissae.

If the abscissae agree, the ordinates are equal or opposite.
Then $P-R$ is $O$ or $-2R$, whose height is at most
$\max(1,H(2R))H(P)^2$. This includes a point of order two.
If $P=O$, the needed inequality is $H(R)\le T_R$;
if $R=O$, it is $H(P)\le H(P)^2$. These follow from the definitions.
All exceptional branches have now been covered.
\end{proof}

\subsection{Unconditional finite division by two}
\label{subsec:halves}

Producing the preimages of a given rational point under doubling
is much easier than finding an unknown representative of every quotient
class. We make this distinction algorithmic.

\begin{lemma}[A complete finite rational-root scan]\label{lem:ratroots}
The rational roots of any nonzero integer polynomial of fixed degree can
be listed by a primitive recursive algorithm.
\end{lemma}
\begin{proof}
Record zero exactly when the constant coefficient vanishes, then strip
the maximal power of $X$, using a scan bounded by the degree.
If a nonconstant polynomial
$c_dX^d+\cdots+c_0$ remains with $c_dc_0\ne0$, a nonzero rational
root $m/n$ in lowest terms, $n>0$, satisfies
\[
 c_dm^d+c_{d-1}m^{d-1}n+\cdots+c_0n^d=0.
 \tag{4.24}
\]
Reducing the equation modulo $m$ gives $m\mid c_0n^d$.
Since $\gcd(|m|,n^d)=1$, $|m|\mid|c_0|$.
Similarly $n\mid|c_d|$. Scan
$-|c_0|\le m\le|c_0|$ and $1\le n\le|c_d|$, test coprimality,
and test (4.24). This is complete and sound.
If stripping leaves a constant, there are no nonzero roots.
The division tests and finite scans are covered by
Lemmas~\ref{lem:gcd} and \ref{lem:scan}.
\end{proof}

\begin{theorem}[Primitive recursive rational halving]\label{thm:halves}
There is a fixed primitive recursive function $\Halves_E(T)$ that lists
exactly the rational points $Q$ with $2Q=T$. In particular rational
divisibility of a given point by two is decidable.
\end{theorem}
\begin{proof}
First take an affine target whose reduced abscissa is $u/v$, $v>0$.
By \eqref{eq:dupmap}, the abscissa $X=x(Q)$ of a preimage satisfies
\[
\begin{split}
 p_T(X)={}&vX^4-4uX^3-2AvX^2\\
          &-(8Bv+4Au)X+(A^2v-4Bu)=0.
\end{split}
\tag{4.25}\label{eq:halvingquartic}
\]
Its leading coefficient is $v>0$, so it is nonzero.
List all its rational roots by Lemma~\ref{lem:ratroots}.
For each root $x$, calculate the reduced rational
$z=x^3+Ax+B$. If $z<0$, discard it. Otherwise test the numerator
and denominator of $z$ for exact squares. If they are squares, test
both points $(x,\sqrt z)$ and $(x,-\sqrt z)$ by the full exact equation
$2Q=T$. Keep exactly the successful points, removing duplicates if desired.

Soundness follows from that final check. For completeness, a preimage of
an affine target is affine and is not of order two; its abscissa satisfies
\eqref{eq:halvingquartic}, and its ordinate is one of the two rational
square roots just tested. Thus it is retained.

For $T=O$, include $O$ and the points $(x,0)$ where $x$ is a rational
root of $x^3+Ax+B$. These are exactly the zero-denominator branches
of doubling. The cubic is monic, so the rational-root scan is again
bounded. Every step has explicit input-dependent finite bounds.
There is no search for points of unbounded height.
\end{proof}

An empty output is a genuine decision about that particular
point. It is not a decision of the universal assertion
$\Cover_E(R)$. Pointwise decidability of a quotient equivalence relation
does not by itself supply a complete list of its classes.

\section{The supplied-cover constructive elliptic Mordell--Weil theorem}
\label{sec:cover-thm}

Let $R=(R_0,\ldots,R_{r-1})$ be a nonempty list of rational points of
$E$, and suppose \eqref{eq:cover}. Adding $O$ to the list is always
permitted. Compute $L$ from \eqref{eq:L} and $T_{R_i}$ from
\eqref{eq:TR}, then put
\[
 T=\max_{i<r}T_{R_i},\quad
 D=LT,\quad d=\htwo D,\quad C=d+2.
 \tag{5.1}\label{eq:cutoff}
\]
These constants depend only on the given equation and representative list.

\begin{lemma}[A natural-number contraction]\label{lem:bitcontract}
Whenever $P=2Q+R_i$, one has
\[
          H(Q)^4\le D H(P)^2,\qquad
          2h(Q)\le h(P)+C.
 \tag{5.2}\label{eq:bitcontract}
\]
\end{lemma}
\begin{proof}
By \eqref{eq:inverse-translation}, $2Q=P-R_i$.
Combine \eqref{eq:dupheight} and \eqref{eq:translation}.
For the second inequality, if $h(Q)=0$ there is nothing to prove.
Otherwise the definition of $h$ gives
$2^{h(Q)-1}<H(Q)$, while $D\le2^d$ and $H(P)\le2^{h(P)}$.
Comparing powers of two yields
\[
 4(h(Q)-1)<d+2h(P),
 \qquad
 4h(Q)\le d+2h(P)+3\le2h(P)+2d+4.
 \tag{5.3}
\]
Dividing this last natural-number inequality by two proves the result.
This proof can be written entirely with integer multiplication, powers and
order; no logarithm laws or real analysis are being imported.
\end{proof}

\begin{theorem}[Constructive Mordell--Weil from a supplied cover]
\label{thm:elliptic-cover}
For every $A,B\in\Z$ with $4A^3+27B^2\ne0$, and every finite valid
list $R$ satisfying $\Cover_E(R)$, the explicitly computed list
\[
                G_E(R)=R\mathbin{+\!\!+}\mathcal L_E(2^C)
 \tag{5.4}\label{eq:genlist}
\]
generates $E(\Q)$ by chord-and-tangent addition and doubling.
A fixed primitive recursive algorithm returns, for each rational point
$P$, a finite word over this list evaluating to $P$.
Its arithmetic correctness under the stated covering hypothesis has an
iPRA proof by the displayed quantifier-free invariants.

In an abelian-group presentation with the group laws proved, the same
algorithm returns nonnegative integer coefficients. With
$e=h(P)\dotminus C$ and $t=\bl(e)$, every coefficient is at most $2^t$.
The size of the list is bounded by
\[
 |R|+1+(2^{C+1}+1)2^C(2S\,2^{2C}+1).
 \tag{5.5}
\]
\end{theorem}
\begin{proof}
For an input $P$, scan $i=0,\ldots,r-1$, calculate $P-R_i$ and
$\Halves_E(P-R_i)$, and select the first point found.
This is a fixed primitive recursive operation, totalized with a failure
flag if every list is empty. By $\Cover_E(R)$ there exists a correct pair
$i,Q$; Theorem~\ref{thm:halves} proves the corresponding finite list is
nonempty. Thus the failure branch is never taken on a valid input under
the covering hypothesis. Equation \eqref{eq:inverse-translation}, in
its reverse form, gives $P=2Q+R_i$.
Lemma~\ref{lem:bitcontract} proves the contraction for whichever witness
the scan selects.

Lemma~\ref{lem:ell-enum} supplies the entire base region.
Apply Theorem~\ref{thm:abstract}, or its operation-only word version.
The bound (5.5) is (4.9) with $M=2^C$, plus the representative list.

The only unbounded existential elimination in constructing the selector's
proof comes from the hypothesis $\Cover_E(R)$ at the current point.
It proves the success of a specified finite program; it does not
define a choice function by arbitrary witness selection. All induction
formulas are the Boolean scan invariant or the forward state invariant of
Section~\ref{sec:descent}.
\end{proof}

\subsection{An alternative selector using only bounded-height enumeration}

Even without the rational-root optimization, \eqref{eq:bitcontract}
implies $h(Q)\le h(P)+C$. One may scan all $i<r$ and all
$Q\in\mathcal L_E(2^{h(P)+C})$, testing $2Q+R_i=P$.
Completeness of the enumeration and the same covering hypothesis prove
success. This route is larger computationally but shows directly that
the height bounds alone remove the need for a separately supplied selector.
The native Lean finite-search module formalizes the selection mechanism.

\subsection{A small explicit duplication certificate}

For $E_1:y^2=x^3-x$,
$F=(X^2+Z^2)^2$, $G=4XZ(X^2-Z^2)$. The following two identities
can be checked by expansion:
\[
\begin{split}
 (16X^3-12XZ^2)F+(-5X^2Z-3Z^3)G&=16X^7,\\
 (-12X^2Z+16Z^3)F+(3X^3+5XZ^2)G&=16Z^7.
\end{split}
\tag{5.6}
\]
They give $L=36$. With the four points later proved to be all of $E_1(\Q)$,
$T=16$ and $C=\htwo(576)+2=12$. This is a deliberately loose general
bound: the unconditional proof below gives the exact height bound $H=1$.
The example illustrates that the general constants are explicit, rather
than asserting that their large enumeration is a practical way to solve
this curve.

\subsection{The abstract version for number fields and abelian varieties}

For a fixed number field one may encode $K=\Q[T]/(f)$ by coefficient
vectors of degree less than $\deg f$, with $f$ a supplied irreducible
polynomial. Reduction and field operations then have finite algebraic
implementations. A projective variety can be represented by finitely many
coordinate charts and polynomial membership tests, with fixed formulas
and guards for its group operations. To apply
Theorem~\ref{thm:abstract} to $A(K)$ one must additionally provide:
a natural-valued primitive recursive height; complete finite
bounded-height enumerations; an actual covering list modulo multiplication;
and a primitive recursive decomposition with the proved contraction.

The same arithmetic theorem then produces generators. This is a
relative interface, not a completed construction of heights,
number-field geometry, weak descent, or representatives for arbitrary
$A,K$. Merely supplying defining equations for an abelian variety does
not discharge those obligations. Section~\ref{sec:unconditional} supplies
actual abelian varieties of every finite dimension for which all the
required rational-point data are settled without those extra hypotheses.

\section{An unconditional Mordell--Weil family in iPRA}
\label{sec:unconditional}

We now discharge all arithmetic hypotheses for a concrete elliptic curve,
and identify its group rather than merely bounding its rank.
The proof does not use the supplied-cover theorem, a search cutoff,
Selmer groups, or a theorem about all elliptic curves.
Its engine is a completely bounded version of Fermat's square-area descent.
Related arguments were present in the inspected ProveIt Diophantine and
Faltings packages \cite{proveit-geometry,proveit-fermat};
the complete argument needed here is included.

\subsection{Primitive square-area triangles}

\begin{definition}
A square-area triangle is a positive integer tuple $(a,b,c,s)$ satisfying
\[
       a^2+b^2=c^2,\qquad ab=2s^2.
 \tag{6.1}\label{eq:triangle}
\]
It is primitive if $\gcd(a,b)=1$.
\end{definition}

\begin{lemma}[Normalization]\label{lem:normal-triangle}
Every square-area triangle yields a primitive square-area triangle
$(A,B,C,S_0)$ with $C\le c$, by a specified primitive recursive operation.
\end{lemma}
\begin{proof}
Let $g=\gcd(a,b)>0$ and put $A=a/g,B=b/g$.
Then $g^2\mid c^2$, so Lemma~\ref{lem:squares} gives $g\mid c$.
Set $C=c/g$. The pair $A,B$ is coprime and
$A^2+B^2=C^2$.

Exactly one of $A,B$ is even. They cannot both be even by coprimality,
and cannot both be odd because the only square residues modulo $4$ are
$0,1$, whereas two odd squares sum to $2$. Therefore $K=AB/2$ is an
integer. From $ab=2s^2$ we get $s^2=g^2K$.
Again Lemma~\ref{lem:squares} gives $g\mid s$. Put $S_0=s/g$,
cancel $g^2$, and obtain $AB=2S_0^2$.
All four entries are positive and $C=c/g\le c$.
\end{proof}

\begin{lemma}[Bounded Pythagorean parametrization]\label{lem:pyth}
For a primitive positive triangle $A^2+B^2=C^2$, exchange the legs
if necessary so that $A$ is even. There are computable integers $m>n>0$
with
\[
 A=2mn,\qquad B=m^2-n^2,\qquad C=m^2+n^2,
 \tag{6.2}\label{eq:pyth}
\]
where $\gcd(m,n)=1$ and $m,n$ have opposite parity.
The required square-root scans have bound $C$.
\end{lemma}
\begin{proof}
The parity argument makes $B,C$ odd. The pair $B,C$ is coprime:
a common divisor $d$ would satisfy $d^2\mid C^2-B^2=A^2$,
whence $d\mid A$, so $d=1$. Since $C>B>0$, the numbers
\[
              U=(C+B)/2,\qquad V=(C-B)/2
\]
are positive, $U>V$, and $UV=(A/2)^2$.
A common divisor of $U,V$ divides $C$ and $B$, so they are coprime.
Square splitting gives $U=m^2,V=n^2,A/2=mn$.
Their sum and difference give \eqref{eq:pyth}.

A common divisor of $m,n$ would divide $U,V$; hence they are coprime.
Their squares sum to the odd number $C$, so they have opposite parity.
The nonnegative roots are unique, since
$t^2-r^2=(t-r)(t+r)>0$ when $0\le r<t$.
Both radicands are at most $C$, and their roots can therefore be found
by scans of $0,\ldots,C$. Every division by two has an even numerator.
\end{proof}

\subsection{The explicit descending transformation}

\begin{lemma}[A strictly smaller square-area triangle]\label{lem:fermat-step}
A primitive square-area triangle $(A,B,C,S_0)$ gives a primitive
square-area triangle $(A',B',C',S'_0)$ with $0<C'<C$.
The transformation is a fixed total primitive recursive function after
assigning a default output to invalid inputs.
\end{lemma}
\begin{proof}
Use the preceding parametrization. The area equation becomes
\[
             mn(m+n)(m-n)=S_0^2.
 \tag{6.3}\label{eq:fourfactors}
\]
The four factors are positive and pairwise coprime.
For pairs involving $m$ or $n$, subtract multiples to reduce a common
divisor to one of $m,n$. A common divisor of $m+n,m-n$ divides
$2m,2n$, and hence divides $2$ by a Bezout equation for $m,n$.
Both $m+n,m-n$ are odd, so this divisor is $1$.

Repeated square splitting now supplies positive integers $r,s,t,u$ with
\[
        m=r^2,\quad n=s^2,\quad m+n=t^2,\quad m-n=u^2.
 \tag{6.4}\label{eq:fourroots}
\]
In particular $t>u>0$, and $t,u$ are odd and coprime.
The integer $s$ is even. To prove this, if $m$ were even and $n$ odd,
their square forms would give $m\equiv0,n\equiv1\pmod4$,
so $m-n=u^2\equiv3\pmod4$, a contradiction.
Thus $n$ is even, and its square root $s$ is even.

Define
\[
        A'=(t+u)/2,\quad B'=(t-u)/2,\quad C'=r,\quad S'_0=s/2.
 \tag{6.5}\label{eq:fermat-map}
\]
These are positive integers. Expanding gives
\[
 (A')^2+(B')^2=(t^2+u^2)/2=m=r^2,
 \qquad
 A'B'=(t^2-u^2)/4=n/2=2(s/2)^2.
 \tag{6.6}
\]
A common divisor of $A',B'$ divides their sum $t$ and difference $u$,
so the new triangle is primitive. Finally
\[
                 0<C'=r\le r^2=m<m^2+n^2=C.
 \tag{6.7}
\]

To specify a total program, first test the primitive triangle conditions.
On failure return $(0,0,0,0)$ with a failure bit. Otherwise order the legs,
find $m,n$ as in Lemma~\ref{lem:pyth}, find the four roots in
\eqref{eq:fourroots}, and return \eqref{eq:fermat-map}.
Every root scan may use the original $C$:
$m,n,m-n,m+n\le m^2+n^2=C$.
The proof establishes that every demanded root exists and every exact
division succeeds in the valid branch. An invalid root search also has a
specified default output.
\end{proof}

\begin{theorem}[Square-area impossibility in iPRA]\label{thm:no-square-area}
There are no positive integers satisfying \eqref{eq:triangle}.
\end{theorem}
\begin{proof}
Suppose a solution is given and normalize it using
Lemma~\ref{lem:normal-triangle}. Let $q$ code the resulting primitive
triangle, $\Good(q)$ be its primitive recursive Boolean validity
predicate, and $\rho(q)$ its hypotenuse. Let $D$ be the total map just
constructed, and define $q_0=q,q_{j+1}=D(q_j)$ by primitive recursion.
Under the assumption $\Good(q)=1$, quantifier-free induction proves
\[
           \Good(q_j)=1\quad\wedge\quad \rho(q_j)+j\le\rho(q).
 \tag{6.8}\label{eq:fermat-invariant}
\]
The base is immediate. The step follows from
$\Good(D(q_j))=1$ and $\rho(D(q_j))+1\le\rho(q_j)$.
All relations in the invariant are Boolean tests or equations of the
fixed coded functions.

Instantiate this invariant at the explicit term $j=\rho(q)$.
It implies $\rho(q_j)=0$, while validity requires $\rho(q_j)>0$.
This contradiction discharges the assumed solution.
The proof runs for exactly the initial hypotenuse; it chooses neither a
least counterexample nor an infinite descending sequence.
Proving the negation of a solution by deriving falsity from that solution
is intuitionistically valid.
\end{proof}

\subsection{All rational points of \texorpdfstring{$y^2=x^3-x$}{y2=x3-x}}

\begin{theorem}[Exact rational point classification]\label{thm:e1}
In the rational codes defined above, iPRA proves
\[
 E_1(\Q)=\{O,(0,0),(1,0),(-1,0)\},
 \qquad E_1:\ y^2=x^3-x.
 \tag{6.9}\label{eq:e1points}
\]
In particular every rational point has $x$-height $1$.
\end{theorem}
\begin{proof}
The projective point at infinity is $O$, so consider an affine point
$(x,y)$. Rational equality is decidable. In the branch $y\ne0$,
the equation excludes $x=0,1,-1$. Define the positive rational numbers
\[
 \alpha=\left|\frac{x^2-1}{y}\right|,
 \qquad
 \beta=\left|\frac{2x}{y}\right|,
 \qquad
 \gamma=\left|\frac{x^2+1}{y}\right|.
 \tag{6.10}\label{eq:triangle-map}
\]
They are strictly positive by the exclusions and $x^2+1>0$.
The polynomial identity
$(x^2-1)^2+4x^2=(x^2+1)^2$ gives
$\alpha^2+\beta^2=\gamma^2$. The curve equation gives
\[
 \frac{\alpha\beta}{2}
   =\left|\frac{x(x^2-1)}{y^2}\right|=1.
 \tag{6.11}
\]
Let $d>0$ be the product of the positive denominators of
$\alpha,\beta,\gamma$. Then
$a=d\alpha,b=d\beta,c=d\gamma$ are positive integers and
$a^2+b^2=c^2$, $ab=2d^2$. This contradicts
Theorem~\ref{thm:no-square-area}, with $s=d$.

Thus the decidable case split on $y=0$ leaves only $y=0$.
Then $x(x-1)(x+1)=0$. The integral-domain property of the rational
arithmetic, or finite case distinction on these three factors,
gives $x=0,1,-1$. Direct substitution verifies all three points.
The list is pairwise distinct, and their heights, including $O$, are $1$.

Every operation in this reduction is normalization, exact rational
arithmetic, sign comparison, or the bounded triangle procedure.
The nonexistence proof was already given with its explicit induction
formula. No empirical point enumeration is used to conclude completeness.
\end{proof}

\subsection{The group, generators, relations and coordinate function}

Write $T=(0,0)$, $U=(1,0)$, $V=(-1,0)$. Each has negative equal to
itself, so its double is $O$. The line through any two distinct affine
points in this list is $y=0$. Formula \eqref{eq:add} shows that the sum
is the remaining affine point. Thus the complete table is
\[
\begin{array}{c|cccc}
 +&O&T&U&V\\ \hline
 O&O&T&U&V\\
 T&T&O&V&U\\
 U&U&V&O&T\\
 V&V&U&T&O
\end{array}
\tag{6.12}
\]
Identify these labels respectively with
$(0,0),(1,0),(0,1),(1,1)\in\{0,1\}^2$.
The table is componentwise addition modulo two. This proves all group
laws, the isomorphism
\[
                     E_1(\Q)\cong(\Z/2\Z)^2,
 \tag{6.13}
\]
and the presentation with generators $T,U$ and relations
$2T=2U=0$, $T+U=U+T$. There is no additional relation between the two
bits: their four sums are the four distinct points.

The coordinate algorithm tests equality with the four listed points
and returns the associated pair of bits. By Theorem~\ref{thm:e1} the
successful branch always exists for a valid point; finite search supplies
the witnesses. It therefore proves a genuine constructive finite-generation
statement, including exact coordinates and relations. The finite table
is separately checked in \code{lean/E1FiniteTable.lean}; the arithmetic
surjectivity onto that table is the preceding written proof.

\subsection{Uniform rational scalings and finite products}

\begin{theorem}[An unconditional family in every finite dimension]
\label{thm:products}
For every rational $u\ne0$ let $E_u:y^2=x^3-u^4x$. Then iPRA proves
\[
 E_u(\Q)=\{O,(0,0),(u^2,0),(-u^2,0)\},
 \qquad E_u(\Q)\cong(\Z/2\Z)^2.
 \tag{6.14}
\]
For every finite input list $u_1,\ldots,u_r\in\Q^\times$, the rational
group of the product has the explicit presentation
\[
 \left(\prod_{j=1}^rE_{u_j}\right)(\Q)
       \cong(\Z/2\Z)^{2r}.
 \tag{6.15}
\]
It has exactly $4^r$ points, rank zero, and a primitive recursive
$2r$-bit coordinate map. The $2r$ coordinate-axis points are generators.
For $r=0$ this means the trivial group and an empty generator list.
\end{theorem}
\begin{proof}
The mutually inverse rational maps
\[
 (x,y)\longmapsto(u^2x,u^3y),\qquad
 (X,Y)\longmapsto(X/u^2,Y/u^3)
 \tag{6.16}
\]
fix $O$ and carry the respective curve equations to one another.
Their denominators are nonzero because $u\ne0$.
Theorem~\ref{thm:e1} therefore proves the displayed list for $E_u$.
Its three nonidentity ordinates are zero; applying
\eqref{eq:add} gives the same four-element table. Thus the
coordinate bijection is a group isomorphism, without needing any
separate abstract theorem about coordinate changes.

A rational point of the product is a tuple of points in its factors.
Apply the four-case test independently to every coordinate, obtaining
$2r$ bits. Conversely any $2r$ bits give exactly one such tuple.
Both transformations are bounded list folds, inverse by their
coordinatewise identities. Addition is bitwise modulo two, so the
coordinate-axis tuples generate and the number of points is $2^{2r}=4^r$.
Induction on the finite input list can be expressed by a Boolean test
of all already processed coordinates together with the accumulated
integer count; it is quantifier-free after list coding.

The ordinary geometric description of this product as an abelian variety
of dimension $r$ is background terminology. The iPRA assertion proved
here is the explicit statement about rational tuple codes, their
membership tests, group table and coordinate map; it does not quantify
over schemes.
\end{proof}

The models in this family are rationally isomorphic to powers of the same
elliptic curve. This is an unconditional infinite family of input
presentations, not a claim that arbitrary twists, positive-rank curves,
or arbitrary number fields have been settled.

\section{Lean formalization: completed kernel and exact remaining targets}
\label{sec:formalization}

\subsection{The distinction between three kinds of verification}

There are three separate claims in this package.
First, the paper supplies human-readable arithmetic proofs whose programs
and induction invariants are primitive recursive and quantifier-free.
Second, selected general algorithms and finite-table assertions are proved
in ordinary Lean. Third, one can encode the arithmetic proofs in the
first-order object theory implemented by ProveIt.
The first two are delivered; the full third development is a detailed
implementation plan, not a completed proof term.

An ordinary theorem of Lean about natural numbers is not automatically a
derivation in iPRA. In particular, native Lean permits induction on
propositions with quantifiers, even if the proof never invokes
\code{Classical}. The embedded target must instead construct a term of
\[
 \code{Prf PRAAx}\ \Gamma\ \varphi,
 \qquad\text{written }\Gamma\vdash[\code{PRAAx}]\varphi.
 \tag{7.1}\label{eq:7-derivation}
\]
ProveIt's \code{pra\_iff\_native} translates a derivation that already
exists in one of its embedded calculi. It is not a reflection rule that
turns every ambient Lean theorem into \eqref{eq:7-derivation}.

\subsection{Repository foundations inspected at a fixed revision}

The inspected ProveIt revision is
\begin{center}
\code{27da047166883e5010e115968d3a40f9e91bfb3d}.
\end{center}
The repository was read through the selected GitHub connector. No upstream
files were changed, and the entire repository was not rebuilt.
The following are the exact relevant files under
\path{Logic/PrimitiveRecursiveArithmetic/ConciseIPRA/Lean/ConciseIPRA/}:

\begin{longtable}{@{}p{0.30\textwidth}p{0.64\textwidth}@{}}
\toprule
File & Role in the intended proof\\
\midrule
\endhead
\code{Syntax.lean} & \code{Fm} and the syntactic predicate \code{Fm.QF};
both quantifiers are excluded from the induction formula.\\
\code{Deduction.lean} & Embedded intuitionistic natural deduction,
equality, substitution, weakening and proof composition.\\
\code{PR.lean} & Function-symbol grammar \code{PR}, terms \code{PTm},
axioms \code{PRAAx}, and the derived proof constructor \code{PRA.qfInd}.\\
\code{PRAToScalar.lean} & Compilation of a specified primitive recursive
definition into the alternative iterator language.\\
\code{TeqReflect.lean} & Reflection of quantifier-free formulas into
characteristic terms and the corresponding induction.\\
\code{BiInterpretation.lean} & Translation between the two object theories.\\
\code{NativeBiInterpretation.lean} & Translation between already
constructed embedded derivations, including \code{pra\_iff\_native}.\\
\bottomrule
\end{longtable}

The existing Diophantine-geometry package supplies written bounded-height,
resultant and square-area arguments and a native arithmetic kernel.
Its earlier finite-generation interface assumes an already supplied
primitive recursive reducing selector. The present elliptic development
derives that selector from the actual representative list and its covering
proof. The complete rational-halving algorithm is a further explicit
piece of this derivation. The square-area proof here follows the especially
self-contained version in the repository's Faltings package, with its
exact fixed-fuel invariant. These dependencies and the audit scope are
recorded in \code{metadata/repository\_audit.txt}.

\subsection{The delivered native Lean definitions}

The generic point type is an arbitrary \code{Type}; its laws are supplied
in a structure. The central interface in \code{DescentCore.lean} contains:
\begin{lstlisting}
structure DescentData (alpha : Type u) where
  add : alpha -> alpha -> alpha
  height : alpha -> Nat
  cutoff : Nat
  repCount : Nat
  baseCount : Nat
  rep : Fin repCount -> alpha
  base : Fin baseCount -> alpha
  repIndex : alpha -> Fin repCount
  baseIndex : alpha -> Fin baseCount
  child : alpha -> alpha
  baseComplete :
    forall p, height p <= cutoff -> base (baseIndex p) = p
  decompose :
    forall p, cutoff < height p ->
      add (add (child p) (child p)) (rep (repIndex p)) = p
  contract :
    forall p, cutoff < height p ->
      2 * height (child p) <= height p + cutoff
\end{lstlisting}
This listing transliterates the actual Unicode Lean source for readability.
There is no declaration asserting that every elliptic curve has such a
structure.

The inductive datatype \code{Word r b} has base and representative atoms,
addition, and a compact doubling constructor. \code{Word.eval} interprets
it in the given operation; \code{Word.size} counts constructors;
\code{baseCoeff} and \code{repCoeff} compute its nonnegative coefficients.
The definitions \code{next} and \code{run} freeze inside the base region.
\code{build} recursively constructs the word with numerical fuel, and
\code{coordinates p = build (height p) p}.

The completed theorems include
\[
\begin{split}
 h(\code{run}(n,P))&\le C+(h(P)\dotminus n),\\
 h(\code{run}(h(P),P))&\le C,\\
 \ev(\code{coordinates}(P))&=P,\\
 \operatorname{size}(\code{coordinates}(P))&\le3h(P)+1.
\end{split}
\tag{7.2}
\]
The proof is independent of the group laws. The stronger logarithmic
fuel bound in Theorem~\ref{thm:abstract} is proved in the paper and used
in the arithmetic reference algorithm; it is not advertised as a theorem
already contained in the native kernel.

\code{FiniteSearch.lean} defines a structurally recursive \code{first?},
proves soundness and completeness on an explicit finite list, then supplies
total selection with a default and bounded-natural search.
\code{E1FiniteTable.lean} defines four labels and checks the table,
commutativity, associativity, exponent two, and a two-generator coding.
It does not infer that these labels exhaust the rational curve.
\code{DescentExample.lean} executes the kernel on a simple numerical
instance; \code{Audit.lean} prints the actual axiom dependencies.

The sixth source, \code{EllipticSpec.lean}, supplies concrete native
definitions of integer, pair, rational and point codes; curve membership;
chord addition, negation and doubling; $H$ and bounded-search $h$; the
halving quartic; the covering predicate; finite division candidates; and
generator-word syntax. Its remaining arithmetic obligations are
definitions of propositions, including \code{HalvesCorrectTarget} and
\code{E1ClassificationTarget}. They are not asserted axioms or proved
theorems. Native \code{Rat} is used as an implementation specification;
its standard-library proof fields have dependencies from Lean's fuller
foundations, including classical choice. The first-order arithmetic bridge
must replace these with the explicit numerical coding proofs. This does
not affect the separate descent/search axiom audit.

Compilation used Lean 4.32.0, commit
\code{8c9756b28d64dab099da31a4c09229a9e6a2ef35}.
The main descent and finite-search theorems report only
\code{propext} and \code{Quot.sound}; the finite table's
\code{two\_generators} theorem reports no axioms.
There are no \code{sorry} placeholders or new axiom declarations in the
compiled files, and no reported \code{Classical.choice} or \code{sorryAx}.
This is not a claim that the entire arithmetic paper is formally proved.

\subsection{All data needed for the arithmetic bridge}

Use the concrete encodings of Section~\ref{sec:foundations}.
The intended embedded development has fixed symbols for:

\begin{enumerate}
\item signed integer operations, positive division, gcd with Bezout
coefficients, rational normalization and ordered comparisons;
\item pair/list operations, bounded scans and finite polynomial evaluation;
\item \code{valid(A,B,P)}, \code{neg}, \code{add}, \code{twice},
and the projective abscissa, with formulas \eqref{eq:add};
\item $H,h,\mathcal L_E$, the matrix $M_{A,B}$, its determinant,
adjugate multipliers, normalized $L$, $T_R$ and $C$;
\item rational-root enumeration, rational-square testing and
$\Halves_E$, including the exceptional target $O$;
\item the finite-list selector, stopped descent state, compact word,
coefficient vector, and evaluation;
\item the primitive triangle predicate, normalized triangle,
Pythagorean parameters, square-area map and its iterate;
\item the four-point classification code, scale maps and finite-product
coordinate maps.
\end{enumerate}

These definitions are fully specified by formulas and bounded loops in
the paper. Integer square root may use an exhaustive scan for the object
proof even when Python uses a faster implementation.
For words, code the constructors by a finite tag and the pairing function.
Choose each constructor's code larger than its subword codes.
Evaluation and structural tests can be implemented by course-of-values
primitive recursion, storing the computed values for all smaller codes
in a finite list. Thus the native inductive datatype is not assumed as an
extra first-order sort.

Define the object-language cover formula by
\[
 \forall p\bigl(\operatorname{valid}(A,B,p)=1\to
   \exists i\,\exists q\,
   [\,i<|R|\wedge\operatorname{valid}(A,B,q)=1
    \wedge p=\operatorname{add}(\operatorname{twice}(q),R_i)\,]\bigr).
 \tag{7.3}\label{eq:objectcover}
\]
Bounded list membership in other formulas means equality of a primitive
recursive Boolean fold to $1$. In particular the successful selector
predicate and all induction invariants are quantifier-free.

The main elliptic target is a derivation under
\eqref{eq:objectcover}, nonsingularity and validity of the given list.
Abbreviate the computed terms by
$W=\operatorname{coordinates}(A,B,R,p)$ and $G=G_E(R)$. The target is
\[
 \operatorname{valid}(A,B,p)=1
 \ \longrightarrow\
 \bigl[\operatorname{wordValid}(W,G)=1\ \wedge\ \ev(G,W)=p\bigr].
 \tag{7.4}
\]
Universal closures over the input codes are introduced after the proof.
For the unconditional target replace the hypotheses by
$\operatorname{valid}(-1,0,p)=1$ and conclude that the four-case classifier
decodes to $p$. This gives the bit-coordinate theorem directly.

\subsection{A dependency-ordered implementation path}

\begin{longtable}{@{}p{0.21\textwidth}p{0.73\textwidth}@{}}
\toprule
Module to implement & Exact mathematical obligation and proof route\\
\midrule
\endhead
\code{Codes} & Encode integers, pairs, lists and rationals; prove projection
identities, normalization idempotence, and reduced-fraction uniqueness.\\
\code{Arithmetic} & Order, divisibility and gcd; prove the Euclidean
coefficient invariant and the three square-splitting lemmas.\\
\code{FiniteLoops} & Compile scans and folds into \code{PR} terms;
prove their Boolean prefix invariants with \code{PRA.qfInd}.\\
\code{EllipticOps} & Define the total guarded chord formulas, closure,
inverse translation and projective duplication. General associativity is
needed only for the coefficient-vector group presentation.\\
\code{PointBoxes} & Prove the square/cube denominator lemma and the
explicit $(a,b,d)$ enumeration bounds.\\
\code{PolyCertificates} & Fixed-size determinant and adjugate identities,
the explicit discriminant polynomial, normalization and coefficient norms.\\
\code{HeightBounds} & Prove image-gcd divisibility, duplication lower
bound, translation upper bound including the exceptional cases, and
binary-height contraction.\\
\code{Halving} & Prove the rational-root divisibility bounds, square test,
quartic completeness and the exact theorem
$Q\in\Halves_E(T)\leftrightarrow Q\in E(\Q)\wedge2Q=T$.\\
\code{CoverSelector} & From the supplied cover, prove success of the
finite scan through $R_i,\Halves_E(P-R_i)$ and prove its contraction.\\
\code{DescentIPRA} & Prove the forward-state invariant by quantifier-free
induction; construct the final word or coefficient vector and its bound.\\
\code{SquareArea} & Normalize triangles, derive Pythagorean parameters,
prove \eqref{eq:fermat-map}, and derive the contradiction at
the initial hypotenuse using \eqref{eq:fermat-invariant}.\\
\code{E1AndProducts} & Clear denominators in \eqref{eq:triangle-map};
prove all points lie in the four-element table; define inverse bit maps,
scale maps and product folds.\\
\bottomrule
\end{longtable}

Each row has explicit premises and a mathematical proof in this paper.
For example, the target named \code{Halves\_correct} is a missing
formal arithmetic lemma, not a theorem already present in the compiled
kernel. Once it and the height lemmas are constructed, a target
\code{Cover\_implies\_search\_success} is proved by existential elimination
from \eqref{eq:objectcover}, finite-scan completeness, and inverse
translation. No unbounded search or choice axiom is required.

\subsection{The critical induction translation}

Native Lean's proof of word correctness inducts on fuel while generalizing
over the point. Its induction proposition therefore has a universal
quantifier. The safe embedded implementation instead fixes the initial
point $P$ and numerically iterates one forward state. At stage $j$, a
coarse invariant is
\[
\begin{gathered}
 \operatorname{valid}(Q_j)=1,\quad k_j\le j,\quad
 h(Q_j)\le C+(h(P)\dotminus j),\\
 P=2^{k_j}Q_j+\sum_l a_{j,l}R_l .
\end{gathered}
\tag{7.5}\label{eq:embedded-invariant}
\]
The symbol for the finite sum, all state entries, and validity are fixed
primitive recursive terms. Hence this is a quantifier-free formula with
$P,R,C$ as parameters. For the operation-only version maintain equality
between $P$ and an accumulated context applied to $Q_j$; the context is
a coded compact word with a hole. Its update substitutes
$(Q_{j+1}\oplus Q_{j+1})\oplus R_i$ for the hole.
One can also retain the entire finite trace and verify its equations by a
Boolean fold.

For the optimized fuel, replace the height line by
$2^j(h(Q_j)\dotminus C)\le h(P)\dotminus C$.
For the Fermat proof use exactly \eqref{eq:fermat-invariant}.
These choices resolve the weak-induction issue explicitly rather than
assuming native Lean induction is an admissible object rule.

\subsection{Completion criteria}

A complete embedded formalization should deliver actual \code{Prf PRAAx}
values for the relative elliptic theorem and unconditional classification,
with the primitive recursive symbols and quantifier-free side conditions
defined in the repository. Rebuild those modules, inspect every added axiom
and placeholder, and only then apply the existing interpretation bridges.
The supplied native files and executable arithmetic are useful intermediate
artifacts; they do not meet this final criterion on their own.

Completing the missing finite algebra and coding is a substantial but
specified formalization project. Removing the representative-cover premise
for all elliptic curves is a different mathematical problem and is not
presented as a routine remaining Lean task.

\section{Reference implementation, checks and remaining limitations}
\label{sec:validation}

\subsection{Exact executable arithmetic}

The Python companion uses integers and \code{fractions.Fraction}.
\code{code/elliptic.py} implements the guarded group operations, integer
height, complete finite point boxes, Sylvester determinants and normalized
adjugate identities, translation constants, and the complete
\code{halves} procedure. \code{division\_by\_representatives} scans a
supplied list; it reports failure if no decomposition is found.
This behavior does not assert that the list covers the entire curve.

\code{code/coordinates.py} implements the full frozen trace and coefficient
extractor with $\bl(h(P)\dotminus C)$ fuel, as well as the complete
generator enumeration and conversion to a coordinate vector.
It shares its extraction engine with the $\Z^r$ tests, which exercise
nontrivial active steps. \code{code/symbolic\_resultant.py} expands the
$8!$ determinant permutations using exact symbolic coefficient dictionaries;
it collects the polynomial
$4096A^6+55296A^3B^2+186624B^4$, exactly
$256(4A^3+27B^2)^2$. This checks a polynomial identity coefficientwise,
rather than testing it only at selected curves.

The polynomial verifier checks the two displayed integral identities and
the multiplier norm. Curve metadata is not itself certified by that check:
the generated $F,G$ must also be the explicit forms belonging to the
chosen $A,B$. The generator constructs them by \eqref{eq:FG}, and the
tests compare these fields. This keeps an identity certificate's actual
meaning precise.

The full point-box and rational-root algorithms are primitive recursive
in the sense of their proved finite bounds. Some Python interfaces raise
an exception on malformed input; the corresponding object-theory
functions use a failure bit and a default output. Faster standard-library
gcd and square-root implementations can be replaced by the displayed
bounded scans without changing the mathematical theorems.

\subsection{What the checks establish}

The exact elliptic checks include:
\begin{itemize}
\item duplication polynomial certificates on five nonsingular models,
with projective-coordinate height comparisons;
\item agreement of two independently implemented bounded-point
enumerations at several small cutoffs;
\item the complete four-point table for $E_1$, and its enumeration through
$x$-height $256$ as a regression check;
\item exact addition, translation inequalities and division round trips
on modest multiples of $(0,1)$ on $y^2=x^3-x+1$;
\item zero-constant quartic branches, empty halving fibers, and all halves
of $O$ on $E_1$;
\item the normalized $E_1$ identities with $R_*=16,L=36,C=12$.
\end{itemize}
The two arithmetic test modules contain fourteen passing tests in total.
They also enumerate the general-theorem generator list for $E_1$ at
$H\le4096$ and recover coordinates for all four points; that list contains
eight entries because representatives and base points are both retained.
The $\Z^r$ fixtures check active transitions, frozen transitions, coefficient
bounds and explicit failure behavior when no decomposition is supplied.
These tests detect implementation errors. The enumeration to $256$ is
not the reason the rational point list is complete, and the sampled
inequalities do not prove their universal versions. Completeness and
the general bounds rest on the proofs in the paper.

The file \path{lean/verification.txt} records the compiler output,
axiom audit and executable example. All six Lean source files compile.
The exact arithmetic check results,
source revision, mathematical audit and file hashes are included in
\code{metadata/} and \code{SHA256SUMS}. The accompanying README gives
the reproduction commands and the exact delivered file inventory.

\subsection{Precise deliverable status}

\begin{longtable}{@{}p{0.53\textwidth}p{0.41\textwidth}@{}}
\toprule
Result & Status\\
\midrule
\endhead
Primitive recursive height descent and coefficient bounds
& Full written proof; native Lean kernel for the coarse-fuel word theorem.\\
General elliptic duplication and translation bounds
& Full written elementary proof; exact arithmetic certificate implementation.\\
Complete rational halving
& Full written proof and exact implementation; embedded proof remains to build.\\
Generation from supplied actual quotient representatives
& Full conditional written theorem; no assumption of a precomputed selector.\\
$E_1(\Q)$, scaled models and finite products
& Unconditional written arithmetic proof; native finite-table verification.\\
Full \code{Prf PRAAx} proof objects
& Not yet constructed; definitions, invariants and dependency plan supplied.\\
Uniform generator algorithm for arbitrary elliptic curves
& Not established; actual representative construction remains unresolved.\\
All abelian varieties over all number fields
& Abstract relative descent interface only; no unrestricted constructive proof.\\
\bottomrule
\end{longtable}

No claim of mathematical novelty for Fermat descent, standard height
methods, or the classical group classification is intended. The main
contribution of the package is a carefully delimited arithmetic
reconstruction: explicit constants and bounded witnesses replace the
ineffective choices in the descent step, and the unresolved global
input remains visible.

\clearpage
\begin{thebibliography}{9}

\bibitem{wikipedia}
Wikipedia contributors, \emph{Mordell--Weil theorem}.
User-specified statement; accessed November 15, 2026.
\url{https://en.wikipedia.org/wiki/Mordell%E2%80%93Weil_theorem}.

\bibitem{milne}
J.\ S.\ Milne, \emph{Elliptic Curves}, second edition,
especially Chapter IV.
Author-hosted text:
\url{https://jmilne.org/math/Books/EC2.pdf}.

\bibitem{sutherland}
Andrew V.\ Sutherland, \emph{18.782 Introduction to Arithmetic Geometry,
Lecture 25}, MIT, December 5, 2013, especially pp.\ 1--2.
\url{https://www.ocw.mit.edu/courses/18-782-introduction-to-arithmetic-geometry-fall-2013/e011c006f6aeb95083197457da90598d_MIT18_782F13_lec25.pdf}.

\bibitem{magma}
Magma Computational Algebra System,
\emph{Mordell--Weil Groups and Descent Methods}, official handbook,
version 2.30-1, introduction.
Accessed November 15, 2026.
\url{https://docs.magma-maths.org/v2.30-1/ArithmeticGeometry/EllipticCurvesOverQAndNumberFields/mordell-weil.html}.

\bibitem{proveit}
Vladimir Reshetnikov, \emph{ProveIt}, revision
\code{27da047166883e5010e115968d3a40f9e91bfb3d}.
Inspected foundation:
\path{Logic/PrimitiveRecursiveArithmetic/ConciseIPRA/Lean/ConciseIPRA/}.
\url{https://github.com/VladimirReshetnikov/ProveIt/tree/27da047166883e5010e115968d3a40f9e91bfb3d/Logic/PrimitiveRecursiveArithmetic/ConciseIPRA/Lean/ConciseIPRA}.

\bibitem{proveit-geometry}
ProveIt, \emph{Diophantine geometry in intuitionistic primitive recursive
arithmetic}, same inspected revision.
Relevant sources: \code{sections/02\_heights.tex},
\code{03\_dynamics.tex}, \code{05\_descent.tex},
\code{07\_generation.tex}, \code{08\_lean.tex},
and \code{lean/DiophantineCore.lean}.
\url{https://github.com/VladimirReshetnikov/ProveIt/tree/27da047166883e5010e115968d3a40f9e91bfb3d/Logic/PrimitiveRecursiveArithmetic/docs/proofs/diophantine-number-theory/diophantine_geometry_ipra}.

\bibitem{proveit-fermat}
ProveIt, \emph{Faltings iPRA package}, same inspected revision,
\code{sections/fermat\_instance.tex}.
The square-area segment was independently inspected.
\url{https://github.com/VladimirReshetnikov/ProveIt/blob/27da047166883e5010e115968d3a40f9e91bfb3d/Logic/PrimitiveRecursiveArithmetic/docs/proofs/diophantine-number-theory/faltings_ipra/sections/fermat_instance.tex}.

\bibitem{mathlib}
Lean community, \emph{Mathlib elliptic-curve documentation},
\code{WeierstrassCurve.Affine.Point}.
Optional ordinary-mathematics integration reference; it is not used as
an object-level iPRA proof dependency.
\url{https://leanprover-community.github.io/mathlib4_docs/Mathlib/AlgebraicGeometry/EllipticCurve/Affine/Point.html}.

\end{thebibliography}

\end{document}
